%!TEX program = xelatex
\documentclass[a4paper,12pt]{article}
\usepackage{fontspec}\defaultfontfeatures{Ligatures=TeX}
\usepackage[a4paper,vmargin={3cm,3.5cm},hmargin={3cm,3cm}]{geometry}
%-----------------------------------------------------------------------------%
\usepackage{settings}
%-----------------------------------------------------------------------------%
%%% Title %%%
\title{G6411 Recitation Notes: Model Selection}
\author{Jesse Chieh Chen\\\href{mailto:cc5229@columbia.edu}{cc5229@columbia.edu}}
\date{\today}
%-----------------------------------------------------------------------------%

\begin{document}

\maketitle

\noindent
Consider candidate regressions indexed by $m$,
each fitted by \gls{ols} on the same $T$ observations.
Let $k_m<T$ be the number of regression coefficients including intercepts.
Define
\begin{align}
    \mathrm{RSS}_m=\sum_{t=1}^T\widehat e_{t,m}^2
	\quad\text{and}\quad
    \widehat\sigma_m^2=\frac{\mathrm{RSS}_m}{T}.
\end{align}
For regressions with an intercept, adjusted $R^2$ is
\begin{align}
    \overline R_m^2
    =1-\frac{\mathrm{RSS}_m/(T-k_m)}
    {\sum_{t=1}^T(Y_t-\overline Y)^2/(T-1)}.
\end{align}
Its denominator is common across models.
Thus maximizing adjusted $R^2$ is equivalent to minimizing
$s_m^2=\mathrm{RSS}_m/(T-k_m)$.
Taking logarithms and multiplying by $T$ does not change the ranking:
\begin{align}
    T\log s_m^2
    &=T\log\widehat\sigma_m^2-T\log(1-k_m/T)\\
    &=T\log\widehat\sigma_m^2+k_m+O(T^{-1}),\label{eq-adjusted-penalty}
\end{align}
where the expansion holds for fixed $k_m$ as $T\to\infty$.
Adjusted $R^2$ therefore imposes approximately a penalty of one per regression coefficient.

\section{Information Criteria}

Let $\ell_{T,m}(\theta_m)$ denote the \emph{average} log-likelihood in model $m$,
and let $d_m$ count all its estimated parameters.
An information criterion takes the form
\begin{align}\label{eq-general-ic}
    \mathrm{IC}_T(m)
    =-2T\ell_{T,m}(\widehat\theta_m)+C_Td_m.
\end{align}
We select the model with the smallest value.
The first term rewards fit, while the second penalizes complexity.

\begin{definition}[\glsentryshort{aic}]\label{def-aic}
	The \gls{aic} is defined as
	\begin{align}
		\mathrm{AIC}(m)&=-2T\ell_{T,m}(\widehat\theta_m)+2d_m.
	\end{align}
	That is, $C_T=2$ in \eqref{eq-general-ic}.
\end{definition}

\begin{definition}[\glsentryshort{bic}]\label{def-bic}
	The \gls{bic} is defined as
	\begin{align}
		\mathrm{BIC}(m)&=-2T\ell_{T,m}(\widehat\theta_m)+d_m\log T.
	\end{align}
	That is, $C_T=\log T$ in \eqref{eq-general-ic}.
\end{definition}

\begin{remark}[Penalties]
	For homoskedastic Gaussian regressions, adjusted $R^2$, \gls{aic}, and \gls{bic}
	differ in their penalties:
	\begin{center}
		\begin{tabular}{lc}
			\toprule
			Criterion & Penalty per regression coefficient\\
			\midrule
			Adjusted $R^2$ & Approximately $1$\\
			\glsentryshort{aic} & $2$\\
			\glsentryshort{bic} & $\log T$\\
			\bottomrule
		\end{tabular}
	\end{center}
\end{remark}

\begin{remark}[Does \glsentryshort{bic} Select a Smaller Model?]
	For $T>\exp(2)$,
	\gls{bic} penalizes each additional parameter more heavily than \gls{aic}.
    On the same candidate set, with the same likelihoods and parameter counts,
    its selected model cannot have more parameters than the \gls{aic} selection,
    though they could select the same model.
\end{remark}

\begin{definition}[Correctly Specified Model]
	Consider a fixed, finite sequence of nested candidate models, indexed by increasing size.
	A model $m$ is \emph{correctly specified} if the true data-generating process is contained in it.
	Assume at least one candidate is correct.
	The smallest correct model $m_0$ is the one with the fewest parameters among all correct models.
	Models below $m_0$ omit relevant terms; those above it include unnecessary ones.
\end{definition}

\begin{definition}[Selection Consistency]
    A selection rule $\widehat m$ is consistent for the smallest correct model if
    \begin{align}
        \Pr(\widehat m=m_0)\to1
        \quad\text{as}\quad
		T\to\infty.
    \end{align}
\end{definition}

\subsection{Underfitting versus Overfitting}

Write $\widehat\ell_m=\ell_{T,m}(\widehat\theta_m)$.
The distinction between omitting a relevant term and adding an irrelevant one is crucial.
In regular, correctly specified likelihood models, the relevant orders are
\begin{align}
    2(\widehat\ell_{m_0}-\widehat\ell_m)&\pto\delta_m>0
    &&\text{if }m<m_0,\label{eq-underfit}\\
    2T(\widehat\ell_m-\widehat\ell_{m_0})&=O_p(1)
    &&\text{if }m>m_0.\label{eq-overfit}
\end{align}
The first follows from a strictly worse population fit for an underfitted model.
The second is the usual likelihood-ratio behavior under redundant restrictions.
Thus the total fit loss from underfitting is of order $T$,
whereas the total fit gain from overfitting is only of order one.

\begin{theorem}[A Penalty That Separates the Two Cases]\label{thm-selection}
    Suppose the candidate set is fixed, finite, and nested, contains a unique smallest
    correct model $m_0$, and satisfies \eqref{eq-underfit} and \eqref{eq-overfit}.
    If
    \begin{align}\label{eq-penalty-rates}
        C_T\to\infty
        \qquad\text{and}\qquad
        \frac{C_T}{T}\to0,
    \end{align}
    then minimizing \eqref{eq-general-ic} consistently selects $m_0$.
\end{theorem}
\begin{proof}[Proof Sketch]
    For an underfitted model, divide the criterion difference by $T$:
    \begin{align*}
        \frac{\mathrm{IC}_T(m)-\mathrm{IC}_T(m_0)}{T}
        =2(\widehat\ell_{m_0}-\widehat\ell_m)
        +\frac{C_T}{T}(d_m-d_{m_0})
        \pto\delta_m>0.
    \end{align*}
    The penalty saving cannot compensate for the loss of fit.
    For an overfitted model,
    \begin{align*}
        \mathrm{IC}_T(m)-\mathrm{IC}_T(m_0)
        =-O_p(1)+C_T(d_m-d_{m_0})\pto+\infty.
    \end{align*}
    The growing penalty dominates the random improvement in fit.
    Since there are finitely many candidates, $m_0$ beats all of them with probability tending to one.
\end{proof}

\begin{remark}[Constant Penalties Retain Unnecessary Parameters]
	Compare $m_0$ with a larger model containing $q$ additional parameters.
	Under the usual likelihood-ratio conditions,
	\begin{align}
		T_{\mathrm{LR}}=2T(\widehat\ell_m-\widehat\ell_{m_0})\dto\chi_q^2.
	\end{align}
	A criterion with $C_T=C<\infty$ prefers the larger model whenever
	$T_{\mathrm{LR}}>Cq$.
	Consequently,
	\begin{align}
		\Pr\{\mathrm{IC}_T(m)<\mathrm{IC}_T(m_0)\}
		\to\Pr(\chi_q^2>Cq)>0.
	\end{align}
	The chance of retaining unnecessary parameters does not vanish.
	This applies to \gls{aic}, with $C=2$, and to adjusted $R^2$,
	whose expansion \eqref{eq-adjusted-penalty} gives $C=1$.
	It assumes that at least one overfitted candidate is available.
\end{remark}

\begin{example}[One Unnecessary Regressor]
    Compare the smallest correct model with a model that adds one unnecessary regressor,
    so the number of additional parameters is $q=1$.
    \Gls{aic} prefers the larger model with limiting probability
    $\Pr(\chi_1^2>2)\approx0.157$.
    For adjusted $R^2$, the probability is $\Pr(\chi_1^2>1)\approx0.317$.
    These probabilities do not go to zero as the sample grows.
    In contrast, the \gls{bic} threshold is $\log T\to\infty$.
\end{example}

\section{So Why Use \glsentryshort{aic}?}

Finding the smallest correct model is not always the goal.
We may instead want the finite model that predicts best,
especially when no finite candidate captures the true conditional mean exactly.
To motivate this idea, we introduce yet another information criterion: \emph{Mallows' $C_p$}.
% The motivation for \gls{aic} is easier to see through Mallows' $C_p$:
% correct the in-sample fit so that it estimates out-of-sample prediction error.

\subsection{Mallows' \texorpdfstring{$C_p$}{Cp}: In-Sample versus Out-of-Sample Prediction}

Treat the regressors as fixed and write $Y=\mu+e$, where $\E e=0$
and $\var(e)=\sigma^2 I_T$.
The true mean vector $\mu$ need not lie in any candidate model.
Let $P_m$ project onto the $k_m$ regressors in model $m$,
so the fitted values are $\widehat Y_m=P_mY$.

Imagine observing new outcomes at the \emph{same regressor values}:
$Y^{\mathrm{out}}=\mu+e^{\mathrm{out}}$,
where $e^{\mathrm{out}}$ is an independent copy of $e$.
What is the expected prediction error for the new outcomes, using the fitted values from the original sample?
Holding the regressors fixed, expectations average over the original sample
and, for out-of-sample prediction, the independent new outcomes:
\begin{align}
    \E\|Y-\widehat Y_m\|^2
    &=\|(I_T-P_m)\mu\|^2+\E\bigl[e\T(I_T-P_m)e\bigr]\notag\\
    &=\|(I_T-P_m)\mu\|^2+\operatorname{tr}\bigl[(I_T-P_m)\E(ee\T)\bigr]\notag\\
    &=T\sigma^2-k_m\sigma^2+\|(I_T-P_m)\mu\|^2,\label{eq-cp-in}\\
    \E\|Y^{\mathrm{out}}-\widehat Y_m\|^2
    &=\|(I_T-P_m)\mu\|^2+\E\|e^{\mathrm{out}}\|^2+\E(e\T P_m e)\notag\\
    &=\|(I_T-P_m)\mu\|^2+\sigma^2\operatorname{tr}(I_T)+\sigma^2\operatorname{tr}(P_m)\notag\\
    &=T\sigma^2+k_m\sigma^2+\|(I_T-P_m)\mu\|^2.\label{eq-cp-out}
\end{align}
% For \eqref{eq-cp-in}, expand $Y-\widehat Y_m=(I_T-P_m)\mu+(I_T-P_m)e$.
% The cross term has expectation zero because $\E e=0$.
% Then use the symmetry and idempotence of $I_T-P_m$,
% $\E(ee\T)=\sigma^2 I_T$, and $\operatorname{tr}(P_m)=k_m$.
% For \eqref{eq-cp-out}, expand $Y^{\mathrm{out}}-\widehat Y_m=(I_T-P_m)\mu+e^{\mathrm{out}}-P_me$.
All cross terms have expectation zero because $e^{\mathrm{out}}$ and $e$ are independent and mean zero;
the same projection and trace identities give the result.

The in-sample fit benefits from fitting noise, producing the term $-k_m\sigma^2$.
That fitted noise does not help predict the new outcomes; instead, it adds $k_m\sigma^2$ to their prediction error.
Subtracting \eqref{eq-cp-in} from \eqref{eq-cp-out} gives the key identity:
\begin{align}\label{eq-cp-correction}
    \E\|Y^{\mathrm{out}}-\widehat Y_m\|^2
    =\E(\mathrm{RSS}_m)+2k_m\sigma^2.
\end{align}
Therefore, we can actually estimate the expected out-of-sample prediction error from the in-sample residual sum of squares,
provided we add the correction term $2k_m\sigma^2$.

\begin{definition}[Mallows' $C_p$]
    Let $\widetilde\sigma^2$ be a common consistent estimate of the error variance, used for every candidate model.
	Mallows' $C_p$ is defined as
    \begin{align}
        C_p(m)=\frac{\mathrm{RSS}_m+2k_m\widetilde\sigma^2}{T}.
    \end{align}
\end{definition}

\begin{remark}[Interpretation of $C_p$]
    With $\sigma^2$ in place of $\widetilde\sigma^2$, \eqref{eq-cp-correction} shows that
    this is an unbiased estimate of average out-of-sample squared prediction error.
    Thus minimizing $C_p$ aims to select the model with the smallest prediction error.
\end{remark}

\begin{remark}[The Usual $C_p$]
    Mallows' $C_p$ is more commonly written as
    \begin{align}\label{eq-cp-usual}
        C_p^{\mathrm{usual}}(m)
        =\frac{\mathrm{RSS}_m}{\widetilde\sigma^2}-T+2k_m
        =\frac{T}{\widetilde\sigma^2}C_p(m)-T.
    \end{align}
    Since $T$ and the positive variance estimate $\widetilde\sigma^2$ are common to all models,
    the two definitions give exactly the same ranking.
    We use the prediction-error scale to make the out-of-sample interpretation clearer.
\end{remark}

\subsection{\glsentryshort{aic} and Prediction}

To put $C_p$ in information-criterion form \eqref{eq-general-ic},
use \eqref{eq-cp-usual} to write
\begin{align}
	\mathrm{IC}_{C_p}(m)
	=\frac{T}{\widetilde\sigma^2}C_p(m)
	=\frac{\mathrm{RSS}_m}{\widetilde\sigma^2}+2k_m.
\end{align}
This has a fit term and a penalty of two per regression coefficient, just like \gls{aic}.
The common positive rescaling does not change the selected model.
For independent Gaussian errors with variance $\sigma^2$ held fixed,
the maximized average log-likelihood is
\begin{align}
	\ell_{T,m}(\widehat\beta_m)
	=-\frac12\log(2\pi\sigma^2)-\frac{\mathrm{RSS}_m}{2T\sigma^2}.
\end{align}
Thus the fit term in \gls{aic} is
\begin{align}
	-2T\ell_{T,m}(\widehat\beta_m)
	=T\log(2\pi\sigma^2)+\frac{\mathrm{RSS}_m}{\sigma^2}.
\end{align}
The log term is common to all models, so removing it does not change the ranking of candidates.
Adding the penalty and using $\sigma^2$
in $\mathrm{IC}_{C_p}$ gives
\begin{align}
	\mathrm{AIC}(m)
	=\mathrm{IC}_{C_p}(m)+\text{constant}.
\end{align}
Therefore, \gls{aic} and $C_p$ select the same model in this case:
both aim to minimize prediction error rather than identify the smallest correct model.

\begin{remark}[Estimating the Variance Separately]
	Estimating the variance separately in each model instead gives, after dropping common terms,
	\begin{align}\label{eq-regression-ic}
		\mathrm{AIC}(m)=T\log\widehat\sigma_m^2+2k_m,
		\qquad \widehat\sigma_m^2=\mathrm{RSS}_m/T.
	\end{align}
	This need not give exactly the same ranking as $C_p$, but the predictive motivation remains:
	the penalty corrects the overly optimistic in-sample fit,
	rather than forcing the probability of selecting unnecessary regressors to zero.
\end{remark}

\begin{remark}[Different Goals]
    \Gls{bic} targets the smallest correct model, whereas \gls{aic} targets predictive fit.
    If every finite model is an approximation, retaining more terms can improve prediction
    even though it increases the number of estimated parameters.
    This explains why selection consistency alone does not settle which criterion to use;
    it does not mean that \gls{aic} always forecasts better in finite samples.
\end{remark}

\begin{remark}[An Information-Theoretic Motivation]
	The connection to $C_p$ illustrates the predictive purpose of \gls{aic} using squared error.
	This is but one of many ways to motivate the criterion.
	\href{https://doi.org/10.1007/978-1-4612-1694-0_15}{Akaike (1973)} motivates the criterion
	through expected Kullback-Leibler divergence from the true distribution to the fitted model.
	Akaike's work helped establish the modern information-theoretic approach to model selection,
	placing predictive performance at its center.
\end{remark}

\vfill
\section*{Acronyms}
\printglossaries

\clearpage
\appendix

\section{Example: Selecting the Order of an Autoregression}

Observe $Y_1,\ldots,Y_N$ and choose a maximum candidate order $p_{\max}$.
For each $p=0,\ldots,p_{\max}$, fit
\begin{align}
	Y_t=c+\alpha_1Y_{t-1}+\cdots+\alpha_pY_{t-p}+e_{t,p}.
\end{align}
Use the \emph{same dependent observations} for every candidate:
\begin{align}
	t=p_{\max}+1,\ldots,N,
	\qquad T=N-p_{\max}.
\end{align}
The first $p_{\max}$ observations supply the lags.
We do not start each regression at its own date $p+1$,
because that would compare likelihoods and residual sums of squares from different samples.
With an intercept, $k_p=p+1$.
Calculate
\begin{align}
	\widehat\sigma_p^2
	&=\frac1T\sum_{t=p_{\max}+1}^N\widehat e_{t,p}^2,\\
	\mathrm{AIC}(p)&=T\log\widehat\sigma_p^2+2(p+1),\\
	\mathrm{BIC}(p)&=T\log\widehat\sigma_p^2+(p+1)\log T,
\end{align}
and choose the minimizing order.
Here the residual variance uses divisor $T$, not $T-p-1$,
because $\mathrm{RSS}_p/T$ is the Gaussian maximum-likelihood estimate of the error variance.
These are based on the Gaussian conditional likelihood, treating the initial observations as given.

\begin{example}
	With $N=200$ and $p_{\max}=10$, we fit every candidate on dates $11,\ldots,200$,
	so $T=190$ even for the intercept-only model.
	An $\mathrm{AR}(2)$ then uses $Y_{10},Y_9$ as regressors for its first dependent observation $Y_{11}$.
\end{example}

\end{document}
